\begin{textsample}{POS Dim 1 – al – Score 42.00 – al/al\_ef\_56.txt}  \label{ex:f1_pos_156}
3.
CONCEPÇÕES ESPECÍFICAS DO COMPONENTE ADOTADAS NO DOCUMENTO Propõe -se um currículo para o Estado de Alagoas, deve -se ter em \textbf{mente} que o estudante, egresso do 5º ano do ensino fundamental, possui dois perfis : um que, teoricamente, não possui educação linguística em língua \textbf{inglesa} porque a rede pública não oferta obrigatoriamente esse componente curricular nos primeiros anos de escolarização e outro de um estudante que poderá vir de um sistema de ensino, a exemplo da rede particular em que o componente já é ofertado desde os Anos Iniciais ou até mesmo na pré-escola.
Esse estudante será, então, o público nessa nova etapa da escolaridade, sendo \textbf{imprescindível} que sejam considerados suas expectativas e limitações nesse contato com esse componente curricular.
Se antes esse estudante dispunha de um(a) educador(a) para todos os seus componentes curriculares durante seu horário escolar, agora, a realidade é bem distinta.
Nessa fase de transição, em que há mudança de um professor único para um professor por componente curricular com características e didáticas próprias, é comum haver frustrações e dúvidas aliadas às ansiedades próprias da idade.
Assim como os demais componentes da educação básica, um currículo para a língua \textbf{inglesa} também deve estar pautado à luz da BNCC, que tem como, entre outros objetivos, a garantia de que todos os estudantes recebam uma formação humana integral que contribua para a construção de uma sociedade justa, democrática e inclusiva e que, ainda de acordo com a visão do citado referencial, “ aprender a língua \textbf{inglesa} propicia a criação de novas formas de engajamento e participação dos estudantes em um mundo social cada vez mais \textbf{globalizado} e plural, em que as fronteiras entre países e interesses pessoais, locais, regionais, nacionais e transnacionais estão cada vez mais difusos e contraditórios ”. ( BNCC, 2017, p. 119 ) Todos concordamos que a língua estrangeira, em especial a língua \textbf{inglesa}, juntamente com os meios tecnológicos, tem um papel fundamental na promoção não somente do acesso ao outro, mas também do acesso que o outro terá em relação a nós.
Ao juntarmos a língua \textbf{inglesa} aos meios tecnológicos de comunicação \textbf{virtual}, podemos acessar jornais, ver \textbf{entrevistas}, assistir a filmes, passear em museus, fazer compras sem precisarmos sair de casa.
Além disso, podemos conhecer pessoas, estabelecer conversas e construir amizades com pessoas de povos distantes ( RECEB-AL, 2015, p. 98 ).
Nesse referencial, considera -se ainda a importância de ter desde os Anos Iniciais ( do 1º ao 5º ) o ensino de língua \textbf{inglesa}, isso porque entendemos ser \textbf{primordial} garantir o acesso de nossas crianças a esse componente e também para assegurar uma equidade de repertório linguístico dos estudantes a fim de que possam dar continuidade aos estudos de língua \textbf{inglesa} nos anos seguintes, a exemplo do que é visto nas demais redes nas quais o componente já é contemplado desde o início de escolarização.
Para essa etapa de ensino, os objetos de conhecimentos, as \textbf{competências} e as habilidades propostas pela BNCC e também contempladas no nosso referencial devem estar pautados nas \textbf{competências} específicas da área de linguagem, bem como, deste componente curricular devidamente contextualizados com as nossas realidades locais.
A construção de novos \textbf{organizadores} curriculares terão como referência o aqui proposto, possibilitando que as redes de ensino com suas escolas e seus professores construam seu documento respeitando e valorizando as características e interesses de seus estudantes.
As \textbf{competências} propostas para a língua \textbf{inglesa} nos Anos Iniciais devem ser trabalhadas de forma lúdica, interativa e dinâmica.
Nesse sentido, a criança aprende melhor um idioma por meio de atividades ricas em mímicas, pantomima, comandos, contação de histórias, peças teatrais, dança, jogos dentre outros com espontaneidade, uma vez que se encontra no período operatório-concreto, entre os 7 e 11 anos de idade ( PIAGET, 1971 ).
A criança faz da palavra um brinquedo para exercitar o pensamento, comunicar -se com colegas, estimulando o desenvolvendo cognitivo e mostrando sua \textbf{capacidade} intelectual para aprender qualquer idioma.
Asher ( 1968 ) corrobora com o pensamento Piagetiano quando propõe o \textbf{Método} TPR ( Total Physical Response ), que se \textbf{baseia} em parte na crença de que movimentos físicos ligados à linguagem ajudem na apropriação ou retenção do conhecimento.
Vale \textbf{lembrar} que outros \textbf{métodos} e \textbf{abordagens} precisam fundamentar a prática de sala de \textbf{aula} quando se trata de ensinar a crianças.
Destacamos, portanto, o \textbf{Método} MAT ( Model, Action, Talk ) de Ritsuko Nakata.
Nele é enfatizado o uso de movimentos rápidos e intensivos, permitindo assim ao estudante adquirir o maior número de habilidades na língua \textbf{inglesa} no menor tempo possível.
Isso se aplica ao estudante em nível iniciante.
Já dentre as \textbf{abordagens}, apontamos quatro delas : 1.
A \textbf{Abordagem} \textbf{Funcional} ( Functional Approach ), que enfatiza o contexto de uso de \textbf{certas} expressões ; 2.
A \textbf{Abordagem} Comunicativa ( Communicative Approach ), cuja base é a crença de que a linguagem produzida em sala de \textbf{aula} deve ser utilizada para expressar pensamentos e sentimentos \textbf{relevantes} ao estudante ; 3.
A \textbf{Abordagem} Áudio-lingual ( Audio-Lingual Approach ), que tem como \textbf{foco} os sons e a sintaxe do idioma ; e 4.
As \textbf{Abordagens} Gramatical/Estrutural ( Grammatical/Structural Approaches ), que analisa a gramática do idioma.
Para a BNCC ( 2017, p. 119 ), “ o estudo da língua \textbf{inglesa} possibilita aos estudantes ampliar \textbf{horizontes} de comunicação e de intercâmbio cultural, científico e acadêmico e, nesse sentido, abre novos percursos de acesso, construção de conhecimentos e participação social ”, ou seja, o idioma se apresenta em seu caráter formativo, numa perspectiva de educação linguística que deve ser consciente e \textbf{crítica} com suas dimensões pedagógicas e políticas intrinsecamente ligadas.
Doravante, essa língua deve ser difundida não mais como uma língua estrangeira, “ corretamente ” falada apenas por nativos como estadunidenses e/ou britânicos, mas sim, como uma língua franca, ou seja, de comunicação internacional e que não pertence somente a esses grupos, mas que vai além dos limites territoriais, uma vez que, diferentes usuários com seus variados repertórios linguísticos e culturais, deixam suas marcas nessa língua franca ao realizarem suas práticas \textbf{discursivas} mundo afora, o que se \textbf{remete} a uma educação linguística voltada para a interculturalidade que se pontua no (re)conhecimento das ( e o respeito às ) diferenças e \textbf{compreensão} do outro enquanto usuário do idioma, superando qualquer tipo de preconceito linguístico.
Vamos, portanto, considerar o que propõe o documento de referência quando aponta o que deve ser dado ao ensino de LE o caráter formativo, cuja aprendizagem do idioma passa a ter uma perspectiva de educação linguística, consciente e \textbf{crítica}.
A ser construído um novo referencial curricular para a língua \textbf{inglesa}, à luz da BNCC para Alagoas, é importante que se pergunte : o que se faz com com uma língua?
A resposta a isso é : Conversa -se, pergunta -se, escreve -se e lê -se textos de vários gêneros, desde escritos de placas de rua até produções mais \textbf{complexas} e longas, analisa -se, opina -se e são procuradas soluções para \textbf{problemas}, etc.
Enfim, atua -se no mundo por meio da língua, utilizando -a diferentemente dentro de contextos variados, de acordo com os interesses e necessidades que são impostas.
Nesse viés, a língua \textbf{inglesa} pode, também, tornar -se uma grande aliada na consolidação de temas e/ou conteúdos abordados pelos professores das demais áreas ou componentes curriculares, sempre que o planejamento pedagógico for construído de forma articulada entre os \textbf{atores} ( professores de vários componentes, coordenadores pedagógicos etc. ) e fundamentado também no Projeto Político-Pedagógico, sempre que necessário, atualizado pela escola.
Por esse motivo, acredita -se que o estudante possa, desde o início, utilizar uma língua estrangeira, em especial o \textbf{inglês}, na produção escrita e oral com o objetivo de buscar informações, discutir processos e \textbf{verificar} resultados.
Isso \textbf{implica} um investimento na formação inicial e continuada do professor para que melhore sua prática sempre fazendo uso da língua nas suas quatro habilidades, sendo referência para que o estudante também as desenvolva, enfim, utilizar a língua para a descoberta dos sentidos.
Participar de um chat, escrever uma mensagem de e-mail, dar informações e perguntar são ações que variam de complexidade de acordo com a situação em que acontecem.
Ao dar informações pessoais para alguém com quem se está iniciando uma conversa, o falante não necessita da mesma quantidade de língua que precisaria para dar as mesmas informações pessoais em uma \textbf{entrevista} de emprego.
Entretanto, em um caso ou noutro, as informações pessoais serão basicamente as mesmas, o que permite que se afirme que, em primeiro lugar, o uso da língua varia em complexidade e forma, não somente em decorrência dos contextos em que ela é utilizada em situações formais e informais, mas também com relação de quem se fala/escreve e/ou para quem se fala/escreve.
Em segundo lugar, a língua \textbf{inglesa} dá acesso a informações que, de outra \textbf{maneira}, talvez fosse acessível, favorecendo o conhecimento interpessoal num fluxo ininterrupto de troca de informações de descobertas que contribuem para que sejam enfrentado os desafios de toda ordem que são diariamente impostos.
A isso denomina -se de processos de \textbf{letramentos}, ou melhor, \textbf{multiletramentos}, concebidas especialmente nas práticas sociais e do mundo \textbf{digital}, na qual se entrelaçam e se \textbf{aproximam} diferentes semioses e linguagens ( verbal, corporal, audiovisual ), criando novas possibilidades de identificação e expressão de ideias, valores e sentimentos, que seja mais preciso.
Consideramos que um currículo de uma língua estrangeira para a Educação Básica deve ser \textbf{baseado} nos usos que fazemos da língua, permitindo ao estudante assumir o papel de protagonista na ampliação e/ou construção de seu próprio conhecimento contando sempre com o auxílio do professor em interações significativas para ele e para aqueles com quem entra em contato, oferecendo a contínua possibilidade do acesso ao outro e do outro em relação a si, propiciando reflexões sobre si mesmos e seus valores.
Deve, ainda, possibilitar a inclusão pelo acesso que essa língua, utilizada em seus vários níveis e contextos, possa oferecer, assim como considerar que o conhecimento da língua \textbf{inglesa} se constrói por meio de seu uso concreto em manifestações orais e escritas, divididas na BNCC pelos eixos \textbf{organizadores} ( Eixo Oralidade, Eixo Leitura, Eixo Escrita, Eixo Conhecimentos linguísticos e Eixo Dimensão \textbf{Intercultural} ) em múltiplas situações possíveis de uso e de ensino e aprendizagem presentes no currículo da língua \textbf{inglesa} nos Anos \textbf{Finais}.
Essa organização se dá por meio de unidades temáticas, objetos de conhecimento e habilidades, distribuídos por ano de escolaridade ( 6º, 7º, 8º e 9º anos ), em um crescente grau de complexidade e consolidação das aprendizagens para facilitar a organização didática dos professores.
Considerando, por fim, os pressupostos da área de linguagem e em articulação com as \textbf{competências} gerais da BNCC ( 2017, p. 244 ) e as \textbf{competências} específicas da área de linguagens, esse componente deve garantir aos estudantes o desenvolvimento de \textbf{competências} específicas a seguir : 1.
Identificar o lugar de si e o do outro em um mundo plurilíngue e multicultural, refletindo, criticamente, sobre como a aprendizagem da língua \textbf{inglesa} contribui para a inserção dos sujeitos no mundo \textbf{globalizado}, inclusive no que concerne ao mundo do trabalho. 2.
Comunicar -se na língua \textbf{inglesa}, por meio do uso variado de linguagens em mídias impressas ou \textbf{digitais}, reconhecendo -a como ferramenta de acesso ao conhecimento, de ampliação das perspectivas e de possibilidades para a \textbf{compreensão} dos valores e interesses de outras culturas e para o exercício do protagonismo social. 3.
Identificar similaridades e diferenças entre a língua \textbf{inglesa} e a língua materna/outras línguas, articulando -as a aspectos sociais, culturais e identitários, em uma relação intrínseca entre língua, cultura e identidade. 4.
Elaborar repertórios linguístico-discursivos da língua \textbf{inglesa}, usados em diferentes países e por grupos sociais distintos dentro de um mesmo país, de modo a reconhecer a diversidade linguística como direito e valorizar os usos heterogêneos, híbridos e multimodais \textbf{emergentes} nas sociedades \textbf{contemporâneas}. 5.
Utilizar novas \textbf{tecnologias}, com novas linguagens e modos de interação, para pesquisar, selecionar, compartilhar, posicionar -se e produzir sentidos em práticas de \textbf{letramento} na língua \textbf{inglesa}, de forma ética, \textbf{crítica} e responsável. 6.
Conhecer diferentes patrimônios culturais, materiais e imateriais, difundidos na língua \textbf{inglesa}, com \textbf{vistas} ao exercício da fruição e da ampliação de perspectivas no contato com diferentes manifestações \textbf{artístico-culturais}.
Na BNCC ( 2017, pp. 246-261 ), nos \textbf{quadros} que apresentam as unidades temáticas, os objetos de conhecimento e as habilidades definidas para cada ano escolar do 6º ao 9º, cada habilidade é identificada por um código alfanumérico como no exemplo seguinte : EF06LI01, em que o primeiro par de letras ( EF ) é usado para indicar a etapa do Ensino Fundamental ; o segundo par de letras ( LI ), o componente curricular, nesse caso a Língua \textbf{Inglesa} ; o primeiro par de números ( 06 ), os anos ( 06 a 09 ) a que se referem, sendo aí no exemplo, o 6º ; e o último par de números ( 01 ) indica a posição da habilidade na numeração sequencial do ano.

% matched lemmas: abordagem, aproximar, artístico-cultural, ator, aula, basear, capacidade, certo, competência, complexo, compreensão, contemporâneo, crítico, digital, discursivo, emergente, entrevista, final, foco, funcional, globalizado, horizonte, implicar, imprescindível, inglês, intercultural, lembrar, letramento, maneira, mente, multiletramento, método, organizador, primordial, problema, quadro, relevante, remeter, tecnologia, verificar, virtual, vista
\end{textsample}
